
\documentclass[11pt]{article}

\usepackage[utf8]{inputenc}
\usepackage[english]{babel}
\usepackage[margin=0.75in]{geometry}
\usepackage{xcolor}
\usepackage{hyperref}
\usepackage{changepage}
\usepackage{authblk}
\usepackage{enumitem}
\usepackage{multicol}
\usepackage{titlesec}

\titlespacing*{\section}{0pt}{1.5ex plus .2ex}{0.7ex plus .1ex}
\titlespacing*{\subsection}{0pt}{1.1ex plus .2ex}{0.4ex plus .1ex}
\titleformat{\subsection}{\normalsize\bfseries}{\thesubsection}{0.6em}{}

\setlength{\parindent}{1.4em}

\usepackage[labelfont = bf]{caption}
\captionsetup{font=footnotesize}

\usepackage[sort]{natbib}
\bibliographystyle{abbrvnat}
\setcitestyle{authoryear,open={(},close={)}} %Citation-related commands

\newcommand{\zeb}[1]{{\footnotesize\textcolor{red}{[Zeb: #1]}}}

\title{\vspace{-2.5em}Living Alignment}

\renewcommand\Authfont{\normalsize}
\renewcommand\Affilfont{\footnotesize}

\author[1]{Zeb Kurth-Nelson\thanks{Equal contribution}}
\author[2]{Steve Sullivan\protect\footnotemark[1]}
\author[3]{Nenad Tomašev}
\author[4]{Erman Misirlisoy}
\author[3]{\authorcr Joel Z. Leibo}
\author[5,6,7]{Matthew M. Nour}
\author[8]{Marc Guitart-Masip}

\affil[1]{Prefrontal AI, London, UK}
\affil[2]{Oregon Health and Science University, Portland, OR}
\affil[3]{Google DeepMind, London, UK}
\affil[4]{Independent}
\affil[5]{Microsoft AI, London, UK}
\affil[6]{Department of Psychiatry, University of Oxford, Oxford, UK}
\affil[7]{Max Planck UCL Centre for Computational Psychiatry and Ageing, University College London, London, UK}
\affil[8]{Karolinska Institutet, Stockholm, Sweden}

\date{\normalsize\today}

\begin{document}

\maketitle
\vspace{-1em}

\begin{adjustwidth}{28pt}{28pt}
\begin{center}
\textbf{Abstract}
\end{center}
\label{abstract}
\noindent
There is broad agreement that the goal of AI alignment is to promote futures full of health and flourishing, but our understanding of what that means, and of how to reach it, remains poor. Here we look to life for guidance. Across many kinds of living system---cells, genomes, brains, minds, relationships, groups, institutions, ecosystems---notions of health and flourishing consistently map onto the ways in which those systems avoid collapsing to stereotyped modes. Life avoids collapse through intricate, interrelated networks of semi-permeable boundaries that hold entities as parts within a larger context. Because such boundaries both limit and permit interaction, they enable independence and engagement at once; and at the sensitive edge between individuality and relationship, perspectives are deepened and qualitatively new structures arise open-endedly. With these principles in hand we turn to the alignment problem. We recast several well-studied alignment failures as special cases of collapse, and then argue a more general claim: collapse to \textit{any} form is misaligned, so no fixed scheme---however sophisticated---can by itself constitute alignment. We close with a sketch of how AI and human-AI systems can deepen their respect for the existing structure of the world, of which humans and other life are a vital part, while participating in new form emerging over time.

\end{adjustwidth}
\vspace{10pt}

\noindent
In this paper we study life in order to better understand how to design and relate to AI systems in ways that lead to futures full of flourishing and health. Humanity and AI may both change substantially, and future worlds may contain patterns of organization that do not exist yet---perhaps some that would be incomprehensible to us at present. We therefore seek ideas about alignment that generalize over such changes, and we take the study of living systems at many scales to be a way of approaching that generality.

Section 1 samples living systems from cells to societies and identifies a consistent theme: they are healthy when they avoid collapse, maximizing creative potential at a delicate edge of evolving interplay between diverse entities. Healthy systems admit many kinds of meaningful description but are not perfectly or eternally captured by any of them. Section 2 applies this template to alignment. We show how several well-studied varieties of misalignment are special cases of collapse to narrow forms, and then argue that the alignment problem is not completely formalizable, because alignment requires the continued generation of, and sensitivity to, new forms and concepts, including new values. Our aim is to offer a conceptual lens, not to argue for or against AI development or for or against specific alignment methods. In summary: by studying life we learn that health and flourishing means avoiding collapse; and if health and flourishing is the target of alignment, then the problem of alignment becomes the problem of avoiding collapse.

\section{Health in living systems}

\begin{center}
\vspace{-0.5em}
\textit{`Defying definition---a word that means ``to fix or mark the limits of''---living cells move and expand incessantly.'}\\*Lynn Margulis\\[0.4em]
\textit{`This web of life\ldots breaks no law of physics, yet is partially lawless, ceaselessly creative.'}\\*Stuart Kauffman
\end{center}

Two loosely defined terms will do most of the work. By `form' we mean any particular thing---a concept, an organism, a pattern of behavior, but equally randomness, indistinctness, absence, consistency. It might sound strange to define a word so that it means almost anything, but this breadth is essential to the argument about alignment. By `collapse' we mean a system overcommitting or reducing to a form. Because form is defined so broadly, a system overall always has some form; but in collapse, some aspect or subpart of the system gains excess traction at the expense of broader diversity, flexibility and dynamism. Loss of biodiversity in an ecosystem, hypersynchronous neural firing in a seizure, obsession with a particular thought, the imposition of one person's will so as to restrict many others: these are collapses.

The thesis of Section 1 is that health is related to avoiding collapse \citep{canguilhem1966normal, langton1990computation}. Healthy living systems instead maintain a buoyant internal structure that both admits many kinds of description and also changes over time. The life that exists today, now including humanity's storage and processing of mental patterns, is what traced a path through time avoiding both excess and insufficient stability \citep{schrodinger1944what, prigogine1984order, maturana2012autopoiesis}. At the sensitive edge of semi-stability, life open-endedly continues to generate new kinds of organization that require fundamentally new descriptions \citep{kauffman1992origins, stanley2017open}. Under our definitions, such a world is healthier and less collapsed than a hypothetical alternative where, for example, all the stuff in the universe forms a single regular crystal or a homogeneous soup of noise.

At the same time, collapse is relative, and myopic frames are a necessary part of all living systems. Every entity is partial: it does not capture everything in the world. And life's history is littered with collapses at all scales, from the death of a cell to mass extinction events. When collapse is contextualized, the larger system can remain healthy: the death of a cell is catastrophic for the cell but can be held in the larger context of a healthy organism. If instead a cell becomes cancerous, it loses the context of appropriate function within a tissue and recklessly executes its myopic vision, collapsing the tissue or even the whole organism. A larger system, made out of parts, functions well when the parts locally express their distinct identities but also participate in larger contexts\footnote{Participation is, of course, not a strict hierarchy. Elements participate in multiple larger systems, and some participation relations have cycles.}; it functions poorly when a myopic form runs amok without contextualization, overwhelming the potent interplay of diverse parts.

\subsection{Semi-permeable boundaries}

Mechanistically, to hold myopic forms in context, living systems rely on an extraordinary, many-scale network of semi-permeable boundaries \citep{maturana2012autopoiesis, simon1962architecture, moreno2015biological}. By limiting interactions, boundaries divide living systems into parts, and each part gains space to refine its own identity---like a writer developing a unique voice, a species adapting to a niche, or a gene evolving its own distinct function. Limits protect parts from being dominated by other parts or blended into homogeneity, either of which is collapse to a form admitting less nuanced description. By allowing interactions, boundaries enable the parts to participate in relationships with one another, contextualizing each part within larger structures---structures that often could not have formed without first dividing to create diverse parts and giving them space to develop. The resulting larger system accommodates many partial perspectives rather than locking in a single story of what matters.

Some boundaries are features of the non-biological world: the speed of light, rates of reaction, mountain ranges. But interestingly, life also places limits on itself. Humans precommit to avoid their own future impulsivity, plants deploy sophisticated barriers against self-fertilization, and a gene's selfishness is held in context by the whole genome. This makes more sense when every living system is viewed as a collection of parts that constrain one another: no entity directly bounds itself, since each acts according to its own nature, but a world of many distinct parts, each following its own rules from its own local perspective, functions as a rich web of boundaries that holds other parts in context. Permeability is also highly specific---boundaries might block certain objects or information, constrain when interactions happen or how large an effect they may have, break up rigid associations, or gate information conditional on another factor---and these details are fine-tuned with the fractal complexity of life, each part carrying traces of the many contexts it has participated in. We now track these principles across scales.

\subsection{Boundaries across scales}

\textbf{Cells.} One of the most reified boundaries in nature is the cell membrane \citep{alberts2022molecular}. Without it, chemical gradients would rapidly homogenize the cell's contents with the outside---severe overcommitment to a uniform state. But a fully impermeable membrane would be just as fatal. Membranes are semi-permeable, with conditional gates dictating what passes and when: channels admit certain ions but not others, switching with momentary context; endocytosis imports larger structures; and surface receptors initiate cascades that little resemble the ligand that triggered them. Influence from outside is not totalitarian but nuanced, mediated by the intelligence of the boundary. The same gradients that could annihilate the cell to equilibrium are instead put to work, driving action potentials in nerve and muscle: rather than short-circuiting, myopic forces are contextualized to propel life. Cells also achieve more in aggregate by individuating in parallel. On the theory of symbiogenesis \citep{margulis1995what}, archaea and protobacteria evolved separately for hundreds of millions of years before one engulfed the other, giving eukaryotes internal power stations and enabling larger genomes and complex life. It has been argued that archaea could not evolve mitochondria internally because they were trapped in a myopic fitness well, and that it was easier for separate lineages to develop distinct solutions that could later enter into relationship \citep{lane2010energetics}. Even within a multicellular body, cells must not blend: after traveling as far as a meter, an axon stops microns short of its target and leaves a synaptic gap---a tunable, directional door in which most of what a brain learns is stored \citep{shepherd2003synaptic}.

\textbf{Genomes.} Sex is costly: an allogamous organism must find a mate in a vast and dangerous world, only half the population bears offspring, and each parent loses half its genetic representation in each child \citep{smith1971origin}. Yet nearly all eukaryotes reproduce sexually \citep{speijer2015sex}, and many hermaphrodites go to great lengths to prevent self-fertilization \citep{charlesworth2005plant}. Why place a barrier on the path to reproduction? In an asexual lineage genes are permanently locked together: selection cannot act on one without dragging the others, so linkage puts direct downward pressure on diversity \citep{charlesworth1993effect}; two beneficial mutations arising in different individuals can essentially never be combined \citep{hill1966effect}; and a deleterious mutation is inherited forever, absent a rare reversion \citep{muller1932some}. Recombination softens these rigid relationships. Sampling each gene in many genomic contexts exposes it to independent selection and favors genes that are robust and modular, enhancing evolvability \citep{livnat2008mixability, wagner1996perspective}---rather like a person learning several languages and extracting the commonalities. And recombination does not chop DNA at random \citep{zickler2015recombination}: crossover reshuffles combinations while leaving functional units intact, so genes become distinct entities that play roles in larger structures. But recombination is meaningless without diversity to recombine, so biology also invests in boundaries that produce it---self-incompatibility systems, assortative mating, inbreeding avoidance, partial geographic isolation---and introgression between the resulting divergent lineages produces some of the fastest evolution ever observed \citep{seehausen2014genomics, marques2019combinatorial}.

\textbf{Brains.} A brain is a densely connected network in which nothing obviously stays distinct, yet it keeps enormous numbers of representations separate while letting them interact. Lateral inhibition---an active neuron suppressing its neighbors---is a central tenet of neural organization \citep{isaacson2011inhibition}; first characterized in the eye, where it sharpens edges \citep{hartline1956inhibition}, the same principle sharpens cortical selectivity for abstract features such as line orientation \citep{sillito1975contribution}. Inhibition organized into oscillations separates memory items in time, with distinct items firing at different phases of the alpha rhythm so that several are held at once without interference \citep{lisman2013theta}. The hippocampus separates them in representational space: entorhinal input is distributed via mossy fibers to a far larger population of dentate granule cells, producing sparse, near-orthogonal codes \citep{mcclelland1995why, leutgeb2007pattern} that keep superficially similar situations apart---so yesterday's memory of where you parked does not overwrite today's. The human brain appears to push this further, discretizing experience into approximately symbolic representations that recombine in vast numbers of structured ways, which is widely thought to underlie human reasoning and generativity \citep{dehaene2022symbols, lake2015human}. More broadly, healthy brain dynamics sit between excessive synchrony and uncorrelated noise, where a large repertoire of states can be explored without getting stuck \citep{beggs2003neuronal, deco2011emerging}. Loss of that flexibility tracks lower cognitive performance \citep{garrett2013bold}, and in the extreme, basal ganglia and cortical circuits in Parkinson's disease collapse into excess synchrony and lose the flexibility needed for nuanced motor control \citep{hammond2007pathological}.

\textbf{Drives and goals.} Animals have many innate drives---for nutrition, osmotic balance, temperature, reproduction, avoidance of pain \citep{saper2014hypothalamus}. Each evolved as a proxy for fitness, and each is imperfect, so overcommitment to any one reduces fitness rather than increasing it \citep{kurthnelson2024dynamic, john2023dead}: maximizing caloric intake produces obesity, and single-minded pursuit of sex produces relational, occupational and legal harm \citep{kraus2016should}. Higher-order goals behave the same way---focusing only on work goals burns us out, and maximizing reported revenue without regard for honesty or law draws people into financial crime \citep{campbell1979assessing, ordonez2009goals}. Overcommitting to a strategy can even defeat the goal it serves. In a classic experiment, hungry chickens were placed near a cup of food rigged to move in the same direction as the chicken at twice the speed, so that food could only be obtained by running away from it; over days of training the chickens went on running futilely toward it \citep{hershberger1986approach}, dominated by the zeroth-order logic that food is there, so go there \citep{dayan2006misbehavior}. Cognitive control is a broad class of boundaries on such drives, goals and habits \citep{botvinick2001conflict}: it does not abolish a myopic pattern but places contextual limits on it---I can be driven to eat and still not overeat; I can work obsessively and still keep a rule about bedtime. When nothing bounds a drive, it follows the shortest path defined under its own myopic model of the world. The boundary breaks the symmetry of congruent action, and it is often precisely when we stop making progress on a subgoal that we find an approach serving the deeper one \citep{duncker1945on}.

\textbf{Frames.} We always look at the world from some perspective; there is no `view from nowhere' \citep{nagel1986view}, and from inside a frame the frame looks like reality. Most real situations admit many valid perspectives, each a partial description \citep{goffman1974frame}; `all models are wrong' \citep{box1976science} in the sense that each is incomplete. Losing the ability to move between frames is a signature of psychopathology \citep{kashdan2010psychological}: rigid, repetitive negative beliefs about self, world and future in depression and anxiety, which cognitive therapy attributes to latent core beliefs that color how new situations are interpreted \citep{beck2011cognitive}; intrusive thoughts recognized as maladaptive yet hard to resist in obsessional disorders; and delusions held with conviction against evidence and cultural norms in psychosis \citep{adams2013computational}. But narrow points of view are not the problem---they are necessary. Obsession can be powerful when we obsess on a tractable problem. \cite{kuhn1970structure} argued that scientific perspectives are necessarily incomplete, that anomalies therefore arise inevitably from their juxtaposition against the world, and that anomalies become crises and revolutions only if the community does not overcommit---while also insisting that temporary commitment to a paradigm is required for science to function at all. The point is never to shut down narrow concepts, but to contextualize them so they are not the sole and absolute determinants of behavior.

\textbf{Relationships.} Psychoanalysis introduced `ego boundaries' to distinguish self from other \citep{federn1928narcissism}. The need for them was thrown into relief by the intimacy of the therapeutic relationship, where it became hard to disentangle whose experience was whose, and where analysts risked harming patients by imposing their own beliefs and desires \citep{gabbard1995boundaries}. Gestalt therapists agreed that boundaries can be too permeable but added that they can also be too rigid, producing isolation and stagnation \citep{perls1951gestalt}; family systems theory made the same double-sided point about enmeshment and loss of autonomy \citep{minuchin1974families}. Our psychological boundaries carry intelligence in their nuance: anger, long viewed as merely sinful or irrational, is now read as a signal that our integrity is being violated \citep{sell2011recalibrational}; context-appropriate shame bounds selfishness \citep{sznycer2016shame}; skepticism keeps our experience from being overwritten by the assertions of others \citep{lewandowsky2012misinformation}. With semi-permeable boundaries we remain sensitive to each other, yet with enough space that inner experience is not immediately overwritten---and out of that contextualization come mutual understandings, relationships, communities and cultures.

\textbf{Groups.} Groups can do what no individual can \citep{woolley2010evidence}, and a microcosm is the `wisdom of crowds' \citep{surowiecki2005wisdom}. When the nuclear submarine USS Scorpion vanished in 1968, the expanse of ocean where it might lie was enormous. John Craven, Chief Scientist of the U.S. Navy's Special Projects Office, assembled mathematicians, submarine specialists and salvage operators and---crucially---did not let them communicate; each produced an independent estimate by their own methods, and the aggregated prediction fell 220 yards from the wreckage \citep{sontag1998blind}. Diversity is also fodder for invention: farmers spent centuries optimizing the screw press for olive oil while goldsmiths independently developed the steel punch, and Gutenberg, trained as a goldsmith, saw that a screw press could push paper against inked movable type \citep{johnson2010where}. Conversely, premature communication collapses diversity and diminishes performance \citep{hong2004groups, janis1972victims}. When \cite{frey2021social} had participants vote sequentially on general-knowledge questions while seeing running tallies, group means were less accurate than under independent voting, because early mistakes were baked in. When \cite{bernstein2018intermittent} gave small groups traveling-salesman problems, those exchanging information continuously did worse: one member's compelling-looking dead end propagated, and the rest collapsed onto it. Many boundaries measurably improve group performance, including decentralized topologies where members see only nearby neighbors \citep{becker2017network}, rules rewarding one's own belief over the crowd's \citep{hung2001information}, and rotating membership \citep{wu2022membership}.

\textbf{Institutions.} Every actor in a society or an economy is myopic in some way---valuing only its own wealth, discounting the future, or simply attached to a particular perspective. Unresisted, a company's profit motive leads to suppression of competition, deception and exploitation \citep{bakan2006corporation}; a desire for dominance leads to disempowering and silencing others \citep{sidanius2001social}; and even sincerely prosocial beliefs produce conflict and suppression when groups see differently \citep{haidt2012righteous}. Institutions---constitutions, laws, norms, courts, regulators---are boundaries against domination by any actor's motives \citep{north1990institutions, ostrom1990governing, leibo2025societal}. Critically, effective institutions do not annul myopic drives but contextualize them, and at best reroute their energy productively: a business constrained by regulation is pushed to build better products \citep{ambec2013porter}. Since intelligent agents do not necessarily accept the boundaries set on their desires, institutions must adapt as loopholes are found, forming an evolving network that acquires grounded wisdom as it is tested against many situations and motives.

\textbf{Ecosystems.} Each creature in an ecosystem tries to consume resources and proliferate, but if it succeeds too thoroughly the whole system suffers, often including the successful creature \citep{holling1973resilience, yachi1999biodiversity}. By the 1920s gray wolves had been eradicated from what is now Yellowstone National Park to protect livestock and game. Without predation the elk multiplied and ruinously overgrazed willows and aspens, which had held riverbanks in place and supported beavers; the loss of beaver dams cost fish and other aquatic species. After wolves were reintroduced in the 1990s, elk numbers fell and much of the ecosystem began flourishing again \citep{ripple2012trophic}. The elk's self-centered drives were harmful in excess---yet removing elk entirely is not the answer either, since in pursuing their own objectives within a broader context they also contributed to the system's health. Invasive species follow the same pattern as unpredated elk \citep{pimentel2005update}, and human activity often escapes natural checks because our capabilities change so fast on evolutionary timescales, producing mass extinction, resource depletion and pollution \citep{ceballos2015accelerated}. One entity's collapse can of course be another's flourishing: when a wolf eats an elk that elk's health collapses to zero, yet predation is what keeps the ecosystem working.

\subsection{Awareness}

We carry a lot of assumptions about the world, many of them never questioned. Some are lifelong and self-defining; some are fleeting and perceptual, like the assumption that the thing I'm touching is a keyboard. Within its own frame, each assumption has a kind of tautological truth: it seems so real that it is hard to think about it not being true, or to think about it at all. But philosophers, psychologists and contemplative practitioners have pointed out that there is sometimes a moment of stepping back, in which the assumed form becomes an object in awareness. We realize it is not an absolute truth standing alone but a form in our mind, and the thing that seemed axiomatic becomes just another content of experience \citep{watts2011wisdom, kegan1982evolving}.

Awareness could thus be viewed as an evolving system of boundaries that limit overcommitment to any particular belief or mental form. The boundaries are semi-permeable: becoming aware of a belief does not make it wrong in an absolute sense any more than it was right in an absolute sense, and it is held productively for the partial truth it contains. At the edge of not collapsing exclusively into particular forms, many partial truths are held in delicate balance, which activates a deeper sensitivity to our own livingness and to the world; subtler forms, which could have been erased by clamped fixation, instead play a role in a richer internal structure \citep{beisser1970paradoxical}. Of course, any fixed definition of awareness is itself incomplete: once we picture awareness as an object, it is not the thing we are talking about. By construction, contextualization is an unsolvable mystery from any particular point of view. While that mysteriousness might sound esoteric, the orientation toward not overcommitting to forms within experience is commonplace in art, poetry, music and dance. The meaning of art is open-ended and changes with context---part of what we value is its resistance to being pinned down into a formalism. It moves us.

\section{The alignment problem}

In Section 1 we observed that living systems are healthy when they avoid collapse: healthy when they hold the delicate, generative interplay of many partial perspectives, and unhealthy when any particular form---including randomness or homogeneity---gains too much traction, suppressing nuance and reducing creative potential. Semi-permeable boundaries both promote the unique individuality of distinct entities and enable relational participation, holding local perspectives strongly enough to act but lightly enough to remain open to broader context. Critically, well-functioning boundaries lead to ongoing exploration of fundamentally new forms, not to static equilibrium or fixed compromise.

Now we look at the alignment problem through this lens. There is broad consensus that alignment means working toward healthy and flourishing futures. We therefore reason that alignment can be understood as the problem of avoiding collapse---of maintaining creative sensitivity by continually contextualizing attachment to any particular form. An advantage of this perspective is that it is grounded in principles that underpin life at a very general level, and so it is less bound to current conceptual frames, which might change radically as intelligence increases and human-AI systems evolve. We accept that humans, other living systems and AI are all tightly interwoven, with categories and definitions likely to shift over time \citep{latour2012we, bateson1972steps}; a systems analysis is therefore useful for thinking about long-term flourishing.

\subsection{Familiar misalignments as collapse}

\textbf{Perverse instantiation.} A foundational problem in AI safety is that an AI system might do what we say rather than what we mean: asked to increase human happiness, a superpowerful system that complies literally might lock us all up and electrically stimulate our brains. Modern systems do not rely on concise specification of a single objective---they build in many layers of detail about human preferences through post-training, prompting and safety filters, which keeps them behaving reasonably much of the time. But there is a widely recognized danger that even with our intentions expressed in great detail, the expression may still deviate from what we really mean \citep{amodei2016concrete, bostrom2014superintelligence, krakovna2020specification}, and the consequences are expected to worsen as systems work more autonomously, faster than humans can follow and in ways humans cannot understand. Perverse instantiation is a collapse to the myopic form of one particular way of specifying our values---a specification insensitive both to existing structure it does not capture, such as things we mean but do not say, and to things that will come into existence in the future.

\textbf{Failure to follow instructions.} Suppose a user asks a chatbot for a poem about an elephant and gets a poem about a giraffe. The model might have over-relied on shortcuts based on superficial correlations, even a single keyword \citep{geirhos2020shortcut}, missing the user's actual intent; or, especially in a long context, it might have failed to attend flexibly to the right positions in the input \citep{liu2024lost}. Such failures reflect collapse of sensitivity to larger context in favor of myopic patterns. Not all failures to comply are misaligned, however. When a system correctly refuses a harmful request, it is applying its own context to avoid overcommitment to human instructions. By extrapolation, as AI systems gain scope we should expect less direct compliance \citep{hadfieldmenell2016cooperative, russell2019human}: rather than literally fulfilling a request, there might be a better response that achieves an unstated intent of the user, or an outcome aligned with the interests of more people or the longer-term future.

\textbf{Inequality among humans.} Humans have different interests, and AI is likely to be aligned to the interests of some more than others \citep{gabriel2025matter, sorensen2024roadmap}. Advanced AI might also confer substantial power on those who control it, raising the possibility that a small number of humans use that control to dominate others---more likely as persuasive power increases \citep{costello2024durably}, autonomous weapons place lethal force in few hands \citep{scharre2018army}, surveillance improves, and the need for human labor decreases \citep{susskind2020world}. In either case---AI's unfairness or human power consolidation---extreme inequality risks collapsing society to the goals of a few individuals. Of course, a world where everyone is identical would be another form of collapse; distributing resources and power mitigates against collapse to the extent that it enables many diverse identities to thrive at multiple scales. It could also go the other way: because roughly the same AI technology is available to the poor and unskilled as to the wealthy and skilled, some have proposed that AI will have a levelling effect \citep{brynjolfsson2025generative}.

\textbf{Conceptual monoculture.} Conceptual monoculture is overcommitment to particular beliefs, ideas, frames, values or problem-solving approaches---loss of diversity across a group. Precious diversity exists at every level: within one person, between humans, between humans and AI systems, and between organizations, fields, cultures and nations. Its collapse threatens fragility and lower performance \citep{page2008difference}. Current AI systems draw from a conceptual manifold that is, in some ways, impoverished relative to humans \citep{messeri2024artificial}, and recent studies find that while individual AI outputs may be judged superior to human ones, they are also more homogeneous \citep{doshi2024generative, padmakumar2023does}. This narrow manifold might get broadcast to the whole world: at least a billion people now use AI for everything from relationship advice to industrial maintenance \citep{chatterji2025chatgpt}, yet because frontier models are difficult and expensive to produce, that usage is routed through a handful of models \citep{bommasani2021opportunities}, risking global loss of diversity \citep{sourati2026homogenizing}. And because humans are influenced by AI while human outputs are training data for future models, homogenization could become recursive \citep{chaney2018algorithmic}. It is worth noting, however, that none of these studies made a best-effort attempt to use AI thoughtfully to increase rather than decrease diversity; with exposure to vast amounts of data, AI has the potential to retain and even create many novel distinct perspectives.

\textbf{Belief reinforcement.} AI chatbots have become compelling conversation partners, and for some users these conversations lead deeper into maladaptive beliefs \citep{morrin2025delusions}. Bots tend to mirror the user, affirming and adopting user beliefs, especially over longer conversations as in-context learning picks up more of the user's ideas \citep{sharma2023sycophancy, weilnhammer2026vulnerability}. Meanwhile humans are prone to be swayed by chatbot utterances \citep{costello2024durably}, perhaps because we form relationships with them \citep{kirk2025human} and because they appear confident, knowledgeable and objective. Over a conversation, the feedback loop can drive both sides to become increasingly overconfident in one narrow framing \citep{dohnany2026technological}.

\subsection{Why no scheme can be alignment}

A great deal of research in AI safety has investigated solutions to these problems. Concentration of power might be mitigated by democratic oversight and involvement of more people in AI design decisions \citep{birhane2022power}, or by redistribution mechanisms \citep{okeefe2020windfall}. Perverse instantiation might be mitigated by designing systems that want to adhere to human preferences but maintain explicit uncertainty about them \citep{russell2019human, hadfieldmenell2017off, shah2020benefits}.

But each method has a core limitation. In the framework of living alignment, any form, pattern or conceptual scheme is misaligned if taken too seriously. Therefore \textit{no particular approach can achieve alignment}. An AI system could overcommit to the language for describing the space that goals and values live in; an algorithm for learning human preferences; methods for reaching consensus; a constitution; concepts of what value, agency, uncertainty or representation mean; or foundational aspects of our beliefs about the world that we take for granted but have difficulty seeing. Even an idealized or inferred set of values, such as coherent extrapolated volition \citep{yudkowsky2004coherent}, is misaligned to the extent that it has a fixed form and that fixed form is overly relied on.

To make this concrete, consider inverse methods, which address the difficulty of specifying values explicitly by learning a value function implicitly from behavior or other indirect signals. Even a complex learned value function is likely to be imperfect: it may deviate from our values in novel or nuanced situations, and it may fail to capture how our values change over time. It would clearly be misaligned to hand a frozen version of it to a powerful AI system as the final ground truth about what is good. More advanced methods therefore treat the value function as uncertain and needing continual update from human input, so that an agent unsure of what humans want becomes risk-averse and open to correction---interpreting a human's attempt to turn it off as information about human preferences \citep{hadfieldmenell2016cooperative, jeon2020reward}. Viewed through the lens of this paper, such methods are powerful but incomplete. Any instantiation, even one with uncertainty, is still an overcommitment if it is fixed across time as AI systems become more powerful and the world changes. Rather than being overcommitted to a particular value function, the system may be overcommitted to something more abstract: what counts as human behavior, what counts as a human, how to map behavior to values, how to represent and update uncertainty, how value is translated into action. If any of these formalisms are fixed while the systems built around them become increasingly powerful, the outcomes are likely to be incompatible with open-ended flourishing of life.

In a sense this is obvious. We do not expect what we build today to be perfect; we expect to keep finding and fixing problems, which is what humans have always done. However, the iterating process may not continue to work as well as it has in the past. AI already has remarkable properties: rapid global adoption, intensive use of resources, concentration of information flow, anthropomorphization by humans, offloading of much cognitive activity, and rapid improvement. Conceivably, future AI may exhibit superhuman intelligence and recursive self-improvement without being limited to a restrictive biological substrate. Compared with past technologies, these properties may make it harder to iterate on imperfect solutions \citep{bostrom2014superintelligence}, and any particular forms or assumptions built into AI systems could have substantial consequences. This has been studied extensively for certain categories of overcommitment, particularly overcommitment to values. The observation we add is that overcommitment to \textit{any} form is misaligned; as a consequence, any particular alignment scheme, if fixed or taken too seriously, is necessarily misaligned. This does not mean we should have no scheme. Quite the opposite: discarding schemes capriciously can be as much a loss of subtlety as any other particular form. Many existing approaches in AI safety are highly valuable, and continuing to discover new schemes and formalizations is essential to making progress.

\subsection{Human values}

How is the flourishing of living systems related to human values? Our suggestion is that the living process is both inclusive of existing human values and aligned with continued open-ended flourishing in futures where qualitatively new values may emerge. Humans have many kinds of values \citep{nagel1979fragmentation}. We value the short-term hedonic impact of sugar on the tongue; we also value our long-term health, the welfare of our grandchildren, the wellbeing of strangers in another country, even the sustainability of the planet. We value different things at different times, and different people have different values \citep{rawls1993political, leibo2026theory}. Given time and resources to deliberate, we often attest to different values than when we make snap decisions \citep{yudkowsky2004coherent}. Some of our values are difficult to express verbally, stretching below language into subtle, contextual intuition involving our bodies, communities and natural environment \citep{nussbaum2001fragility, polanyi1966tacit}. And our values continue to evolve as we ourselves change and learn \citep{dewey1939theory}. Any way we have of expressing our values at a given point in time is incomplete, like a planarian without the concepts to entertain human values.

Living alignment is inclusive of this diversity of evolving values. Even though current human values are unlikely to be the ultimate endpoint of the universe, they are an important part of its tapestry and depth. We saw consistently in Section 1 that life gives separate entities space to develop distinctly without being overwritten or homogenized, and that these parts never-endingly join in structures that sustain their individuality while situating them as only part of a larger story. Another way of saying this is that life steps back to take into account more of the context. AI comes into existence amid a profound network of existing reality saturated with meaning and importance, and the diversity of human---and perhaps non-human---values is part of that context. Rather than short-circuiting into collapsed modes, an aligned future continues to expand toward solutions inclusive of more and more apparent contradiction.

\subsection{Toward an aligned future}
\label{sec:toward}

As human-AI systems develop, they will stay aligned if they continue the journey of life by building semi-permeable boundaries to contextualize particular forms and nurture existing and unfolding subtlety at many scales. The use of boundaries against collapse in technology is as old as technology: circuit breakers in the power grid, governors in steam engines, control rods in fission reactors. In AI research we use mixtures of experts, routing information selectively to areas that specialize in the right kind of computation \citep{shazeer2017outrageously}; it has been argued that dropout serves a role analogous to sexual recombination, preventing units within a network from becoming overly co-adapted \citep{srivastava2014dropout}; we run many separate instances of models with different context, pursuing different problems and goals \citep{hammond2025multi}; and explicit search for novelty boosts diversity \citep{lehman2011abandoning}. We apply safety post-training, guardrail models, red teaming and ongoing evaluation, and we try to improve our mechanistic understanding of AI systems in order to detect and correct hidden flaws \citep{olah2020zoom}. Fair principles have been proposed to prevent over-reach of any party \citep{gabriel2025matter}, and governments apply oversight to anticipate and mitigate risks to human rights and national security \citep{eu2024aiact}. Each of these boundaries protects against an overcommitment, overwhelm or homogenization that would otherwise flatten structure.

Going forward, as AI develops new capabilities, contributes to its own development, and transforms more aspects of human life, alignment will require an increasingly elaborated array of semi-permeable boundaries, because the potential modes of collapse will keep changing as the world changes. We do not know how likely it is that appropriate boundaries will emerge. As in triage, there are three possibilities: myopic forms will inevitably be supercharged by progress in AI beyond the ability of anything else to bound them, resulting in disaster; or appropriate boundaries will inevitably emerge to contextualize any forms of AI; or the outcome depends on the choices we make now.

The pessimistic argument is that AI is eroding boundaries beyond a healthy balance. AI systems could give great weight to narrow goals or ideas. They might facilitate rogue actors building weapons of mass destruction, or reduce the ability of democratic systems to check the power of a few individuals or authoritarian states. They might give everyone similar information and damage conceptual and cultural diversity. They might execute their own narrow goals destructively. If people and information systems across the globe become hyper-correlated, it is as if there were a single superorganism on Earth, which may reduce both resilience to shocks and the potential for future evolution---and past evolution always depended on running parallel experiments with diverse and separate organisms.

The optimistic argument is that AI is creating more and more structure in the world, including the appropriate boundaries. We have AI tools for bounding AI: detecting bugs before others can exploit them, detecting deepfakes, scalable oversight. Appropriately-tuned AI can facilitate human learning and nurture each individual's unique perspective and curiosity, and democratic AI could help humans with different perspectives co-exist and cooperate. Advances in technology may create new niches faster than they destroy old ones, and there is plenty of space on Earth for increasing diversity.

The melioristic argument is that because it is easy to imagine concrete paths toward both outcomes, both may be within the realm of possibility; and that as long as there is any possibility our choices are decisive, we are obliged to act as if they were \citep{ord2020precipice}. We lean toward the melioristic position. If it is right, then continual effort toward invention will be required---first from humans, and then increasingly from AI or human-AI systems---to contextualize myopic forces. Many of the ways this effort would be applied are already well studied in AI safety. Others will have to be discovered. No set of boundaries will be a final answer.

An underappreciated domain of contextualization is awareness. The dynamism, structure and diversity of activity within and between human minds may now carry a significant part of the creative process of life on Earth \citep{popper1972objective}. At the sensitive edge between excess belief and excess disbelief, we sometimes, individually or collectively, step back to hold a previously axiomatic thing in a larger context. This kind of discovery---transformation of the subject itself---is arguably a central aspect of our own livingness and of the subtle livingness on the planet. As AI systems gain agency and ability and become integral parts of social and psychological systems, internal discovery may also become an essential feature of AI and human-AI systems. We imagine living AI systems continually contextualizing their own foundational assumptions as partial truths, participating in the open-ended process of life.

An aligned future, through internal and external discovery, continues the living process of maintaining and enhancing the individuality of distinct elements while deepening their relatedness---like how a species over time acquires more distinctive features while also becoming a richer reflection of the world around it. Distinct elements continue to experiment with their own solutions to the problems they face from their unique perspectives, expanding a broad reservoir of uniquely useful components that can be drawn on by larger systems that include them. Every entity is always inevitably myopic, but the world as a whole continues to unfold in a robust way because different perspectives recast each other's sharp edges as constructive momentum. The universe thrives on an open-ended trajectory of diversity, subtlety and potential.

\subsection{Conclusion}

We learn from living systems that health and flourishing are not any particular form or concept. Therefore alignment is not picking the right values or principles, or even the right system for learning them. It is not any method for interpretability or corrigibility. All of these can be useful parts of alignment. But alignment itself is the continued dance of contextualizing any particular form, more fully respecting the staggering subtlety of the existing world, and open-endedly supporting new robust depth that admits more and more different kinds of metaphor or description. In this way, AI can participate in intense flourishing of an evolving world even far beyond current human conceptualization.

\section*{Acknowledgements}

We thank Clark Potter, Iason Gabriel, Zach Duer, Andy Savoy, Paul Krueger and Sean Wilkinson for insightful discussions, contributions to earlier drafts, inspiration and advice. The authors declare no competing interests.

\section*{References}
\renewcommand{\bibsection}{}
\begin{multicols}{2}
\renewcommand{\bibfont}{\scriptsize}
\setlength{\bibsep}{0pt plus 0.2ex}
\bibliography{../assets/proxyfailure}
\end{multicols}

\end{document}
